
\section{Advanced Activities}

The following activities are intended for strong master's or PhD students.
They emphasize nonlinear modeling, state-space representation, linearization,
exact dynamic response, stability, feedback, modal behavior, and the way
physical design decisions change closed-loop system dynamics. MATLAB, Python,
or an equivalent computational environment may be used where specified.

\subsection*{Problem 1: Nonlinear Pendulum Modeling, Linearization, and Local Feedback}

A torque-actuated pendulum is described by
\begin{equation}
ml^2\ddot{\theta}
+
b\dot{\theta}
+
mgl\sin\theta
=
u,
\end{equation}
where $\theta=0$ is the downward equilibrium and $\theta=\pi$ is the upright
equilibrium. Use
\begin{equation}
m=1.2~\mathrm{kg},
\qquad
l=0.8~\mathrm{m},
\qquad
b=0.12~\mathrm{N\,m\,s},
\qquad
g=9.81~\mathrm{m/s^2},
\end{equation}
with the actuator limit
\begin{equation}
|u(t)|\leq 8~\mathrm{N\,m}.
\end{equation}

\begin{enumerate}
    \item Define
    \begin{equation}
    x_1=\theta,
    \qquad
    x_2=\dot{\theta},
    \end{equation}
    and derive the nonlinear state-space model
    \begin{equation}
    \dot{\mathbf{x}}
    =
    \mathbf{f}(\mathbf{x},u).
    \end{equation}

    \item Determine all equilibrium pairs $(\mathbf{x}_e,u_e)$ for
    \begin{equation}
    -2\pi\leq\theta_e\leq2\pi.
    \end{equation}

    \item Linearize the system about the downward equilibrium and derive
    the matrices $A_d$ and $B_d$.

    \item Linearize the system about the upright equilibrium and derive
    the matrices $A_u$ and $B_u$.

    \item Compute the open-loop eigenvalues of both linearized systems and
    classify their local stability.

    \item Design the upright state-feedback law
    \begin{equation}
    u
    =
    -K
    \begin{bmatrix}
    \theta-\pi\\
    \dot{\theta}
    \end{bmatrix}
    \end{equation}
    so that the linearized closed-loop poles are located at
    \begin{equation}
    -3\pm4j.
    \end{equation}

    \item Derive the resulting closed-loop linearized matrix and verify its
    eigenvalues analytically.

    \item Simulate both the nonlinear and linearized closed-loop systems for
    \begin{equation}
    \theta(0)=\pi+\Delta\theta,
    \qquad
    \dot{\theta}(0)=0,
    \end{equation}
    with
    \begin{equation}
    \Delta\theta
    \in
    \{0.02,\ 0.10,\ 0.30,\ 0.60\}.
    \end{equation}

    \item Determine the largest tested perturbation for which the linearized
    model predicts the nonlinear response with less than $5\%$ normalized
    root-mean-square error over the first five seconds.

    \item Repeat the nonlinear simulations with actuator saturation and
    explain why local pole placement does not guarantee large-angle
    stabilization.
\end{enumerate}


\subsection*{Problem 2: Exact State Transition of a Mass--Spring--Damper System}

Consider
\begin{equation}
m\ddot{x}+c\dot{x}+kx=0,
\end{equation}
with
\begin{equation}
m=2~\mathrm{kg},
\qquad
k=18~\mathrm{N/m},
\qquad
x(0)=1,
\qquad
\dot{x}(0)=0.
\end{equation}

Use the state vector
\begin{equation}
\mathbf{x}
=
\begin{bmatrix}
x\\
\dot{x}
\end{bmatrix}.
\end{equation}

\begin{enumerate}
    \item Derive the state matrix $A(c)$.

    \item Derive the characteristic polynomial and the two eigenvalues as
    functions of $c$.

    \item Determine the damping values corresponding to:
    \begin{enumerate}
        \item underdamped motion,
        \item critical damping,
        \item overdamped motion.
    \end{enumerate}

    \item For the underdamped case, derive a closed-form expression for
    \begin{equation}
    e^{At}
    \end{equation}
    using the Cayley--Hamilton theorem.

    \item Use the result to derive exact expressions for $x(t)$ and
    $\dot{x}(t)$.

    \item Repeat the derivation for the critically damped case, taking care
    of the repeated eigenvalue.

    \item For
    \begin{equation}
    c\in\{2,\ 12,\ 20\}~\mathrm{N\,s/m},
    \end{equation}
    compute:
    \begin{equation}
    \omega_n,
    \qquad
    \zeta,
    \qquad
    \lambda_1,
    \qquad
    \lambda_2.
    \end{equation}

    \item For each damping value, determine analytically or numerically:
    \begin{enumerate}
        \item the first zero-crossing time,
        \item the time of the first displacement extremum after $t=0$,
        \item the $2\%$ settling time.
    \end{enumerate}

    \item Verify the analytical state-transition solution using a
    high-accuracy numerical ODE solver.

    \item Explain why increasing $c$ beyond critical damping can make the
    slowest mode slower even though the system contains more damping.
\end{enumerate}


\subsection*{Problem 3: Allocation of Closed-Loop Dynamics between Plant and Feedback}

Consider
\begin{equation}
m\ddot{x}+c\dot{x}+kx=u,
\end{equation}
with the PD feedback law
\begin{equation}
u=-K_px-K_d\dot{x}.
\end{equation}

Use
\begin{equation}
m=2~\mathrm{kg}.
\end{equation}

The desired closed-loop dynamics are
\begin{equation}
s^2
+
2\zeta_d\omega_d s
+
\omega_d^2
=
0,
\end{equation}
where
\begin{equation}
\zeta_d=0.7,
\qquad
\omega_d=4~\mathrm{rad/s}.
\end{equation}

\begin{enumerate}
    \item Show that all plant--controller combinations producing the desired
    poles must satisfy
    \begin{equation}
    k+K_p=m\omega_d^2,
    \end{equation}
    and
    \begin{equation}
    c+K_d=2m\zeta_d\omega_d.
    \end{equation}

    \item Consider the allocation cost
    \begin{equation}
    C
    =
    \alpha_k k^2
    +
    \alpha_c c^2
    +
    \alpha_p K_p^2
    +
    \alpha_d K_d^2,
    \end{equation}
    with
    \begin{equation}
    \alpha_k=0.04,
    \qquad
    \alpha_c=0.08,
    \qquad
    \alpha_p=0.01,
    \qquad
    \alpha_d=0.02.
    \end{equation}

    Derive the minimum-cost allocation analytically.

    \item Repeat the derivation subject to the passive-safety requirements
    \begin{equation}
    k\geq8~\mathrm{N/m},
    \qquad
    c\geq1.5~\mathrm{N\,s/m}.
    \end{equation}

    \item For the initial condition
    \begin{equation}
    x(0)=1,
    \qquad
    \dot{x}(0)=0,
    \end{equation}
    derive the closed-loop state trajectory.

    \item Derive the control trajectory
    \begin{equation}
    u(t)=-K_px(t)-K_d\dot{x}(t)
    \end{equation}
    for an arbitrary feasible allocation.

    \item Show that different allocations can produce identical
    closed-loop poles and identical state trajectories but different
    control-force histories.

    \item Compute
    \begin{equation}
    E_u
    =
    \int_0^\infty u(t)^2\,dt
    \end{equation}
    for:
    \begin{enumerate}
        \item the unconstrained minimum-cost allocation,
        \item the passive-safety-constrained allocation,
        \item the fully passive allocation $K_p=K_d=0$.
    \end{enumerate}

    \item Suppose the actuator fails at $t=1$ and
    \begin{equation}
    u(t)=0,
    \qquad
    t\geq1.
    \end{equation}
    Simulate and compare the three designs.

    \item Explain why matching closed-loop poles alone is insufficient for
    choosing a plant--controller design.
\end{enumerate}


\subsection*{Problem 4: Full and Reduced Models of a DC Motor}

An armature-controlled DC motor is modeled by
\begin{align}
J\dot{\omega}+b\omega
&=
K_ti-\tau_L,
\\
L\dot{i}+Ri
&=
v-K_e\omega.
\end{align}

Use
\begin{equation}
J=0.02~\mathrm{kg\,m^2},
\qquad
b=0.08~\mathrm{N\,m\,s},
\end{equation}
\begin{equation}
L=0.015~\mathrm{H},
\qquad
R=1.2~\Omega,
\qquad
K_t=K_e=0.25.
\end{equation}

\begin{enumerate}
    \item Using
    \begin{equation}
    \mathbf{x}
    =
    \begin{bmatrix}
    \theta&\omega&i
    \end{bmatrix}^{T},
    \end{equation}
    derive the complete state-space model with voltage $v$ as the control
    input and load torque $\tau_L$ as the disturbance.

    \item Derive the characteristic polynomial of the unforced state matrix.

    \item Explain why one eigenvalue is located at the origin when shaft
    position is included as a state.

    \item Remove the position state and derive the two-state speed-current
    subsystem.

    \item Prove that the speed-current subsystem is asymptotically stable for
    positive
    \begin{equation}
    J,\ b,\ L,\ R,\ K_t,\ K_e.
    \end{equation}

    \item Set $L\dot{i}\approx0$ and derive the reduced first-order speed
    model.

    \item Derive the reduced-model pole and compare it with the slow
    eigenvalue of the full speed-current model.

    \item Simulate the full and reduced models for a $12$-V step input and
    zero load torque. Compare:
    \begin{enumerate}
        \item motor speed,
        \item current,
        \item steady-state speed,
        \item dominant time constant.
    \end{enumerate}

    \item Repeat the comparison for
    \begin{equation}
    L\in\{0.005,\ 0.015,\ 0.08\}~\mathrm{H}.
    \end{equation}

    \item Define a quantitative criterion for deciding when the reduced model
    is acceptable.

    \item Explain how changing rotor inertia $J$, inductance $L$, or torque
    constant $K_t$ changes both plant performance and the control-design
    problem.
\end{enumerate}


\subsection*{Problem 5: Lyapunov Analysis of a Nonlinear Feedback System}

Consider the nonlinear oscillator
\begin{equation}
m\ddot{x}
+
c\dot{x}
+
kx
+
\alpha x^3
=
u,
\end{equation}
where
\begin{equation}
m>0,
\qquad
\alpha>0.
\end{equation}

Apply the PD controller
\begin{equation}
u=-K_px-K_d\dot{x}.
\end{equation}

\begin{enumerate}
    \item Write the nonlinear closed-loop state-space model using
    \begin{equation}
    x_1=x,
    \qquad
    x_2=\dot{x}.
    \end{equation}

    \item Determine all equilibrium points as functions of
    \begin{equation}
    k+K_p.
    \end{equation}

    \item Consider the candidate Lyapunov function
    \begin{equation}
    V(x_1,x_2)
    =
    \frac{1}{2}mx_2^2
    +
    \frac{1}{2}(k+K_p)x_1^2
    +
    \frac{1}{4}\alpha x_1^4.
    \end{equation}

    Derive $\dot{V}$ along closed-loop trajectories.

    \item Prove global asymptotic stability of the origin when
    \begin{equation}
    k+K_p>0,
    \qquad
    c+K_d>0.
    \end{equation}

    \item Explain what changes when
    \begin{equation}
    k+K_p<0.
    \end{equation}
    Determine the additional equilibria and classify them locally.

    \item Replace the ideal controller with
    \begin{equation}
    u
    =
    \operatorname{sat}
    \left(
    -K_px-K_d\dot{x},
    F_{\max}
    \right).
    \end{equation}

    \item Derive the region in the $(x,\dot{x})$ plane in which the actuator
    remains unsaturated:
    \begin{equation}
    |K_px+K_d\dot{x}|
    \leq F_{\max}.
    \end{equation}

    \item Find the largest Lyapunov sublevel set
    \begin{equation}
    \Omega_\rho
    =
    \left\{
    (x,\dot{x}):
    V(x,\dot{x})\leq\rho
    \right\}
    \end{equation}
    that is guaranteed to lie inside the unsaturated region.

    \item Verify the analytical stability conclusions numerically for at
    least four initial conditions inside and outside $\Omega_\rho$.

    \item Explain why local linear eigenvalues cannot establish the global
    stability properties obtained from the Lyapunov analysis.
\end{enumerate}


\subsection*{Problem 6: Modal Dynamics and Actuator Placement in a Two-Mass System}

Two unit masses are connected by springs as described by
\begin{equation}
M\ddot{\mathbf{q}}
+
C\dot{\mathbf{q}}
+
K\mathbf{q}
=
B_fu,
\end{equation}
where
\begin{equation}
M=
\begin{bmatrix}
1&0\\
0&1
\end{bmatrix},
\qquad
K=
k
\begin{bmatrix}
2&-1\\
-1&2
\end{bmatrix},
\end{equation}
and proportional damping is used:
\begin{equation}
C=\beta K.
\end{equation}

Use
\begin{equation}
k=20~\mathrm{N/m},
\qquad
\beta=0.015~\mathrm{s}.
\end{equation}

Consider three actuator placements:
\begin{equation}
B_f^{(1)}
=
\begin{bmatrix}
1\\
0
\end{bmatrix},
\qquad
B_f^{(2)}
=
\begin{bmatrix}
0\\
1
\end{bmatrix},
\qquad
B_f^{(3)}
=
\begin{bmatrix}
1\\
-1
\end{bmatrix}.
\end{equation}

\begin{enumerate}
    \item Derive the four-state model using
    \begin{equation}
    \mathbf{x}
    =
    \begin{bmatrix}
    \mathbf{q}\\
    \dot{\mathbf{q}}
    \end{bmatrix}.
    \end{equation}

    \item Solve the generalized eigenvalue problem
    \begin{equation}
    K\boldsymbol{\phi}_j
    =
    \omega_j^2
    M\boldsymbol{\phi}_j
    \end{equation}
    and compute the natural frequencies and normalized mode shapes.

    \item Transform the mechanical equations to modal coordinates
    \begin{equation}
    \mathbf{q}
    =
    \Phi\boldsymbol{\eta}.
    \end{equation}

    \item For each actuator placement, compute the modal input vector
    \begin{equation}
    \widetilde{B}_f
    =
    \Phi^TB_f.
    \end{equation}

    \item Identify which modes are weakly actuated or completely unactuated
    by each placement.

    \item Apply the initial condition
    \begin{equation}
    \mathbf{q}(0)
    =
    \begin{bmatrix}
    1\\
    -1
    \end{bmatrix},
    \qquad
    \dot{\mathbf{q}}(0)=\mathbf{0},
    \end{equation}
    and determine which mode is initially excited.

    \item For each actuator placement, design full-state feedback
    \begin{equation}
    u=-K_x\mathbf{x}
    \end{equation}
    to move the controllable closed-loop poles as far left as possible
    while satisfying
    \begin{equation}
    \|K_x\|_2\leq 25.
    \end{equation}

    \item Compare the closed-loop responses and control efforts for the three
    actuator placements.

    \item Repeat the modal analysis after changing the second mass to
    \begin{equation}
    m_2=1.4~\mathrm{kg}.
    \end{equation}

    \item Explain how a physical design change can alter mode shapes,
    actuator authority, and the controller that performs best.
\end{enumerate}
